% Generated draft: requires mathematical and source review.
\documentclass[11pt]{article}
\usepackage[margin=25mm]{geometry}
\usepackage{fontspec}
\setmainfont{Latin Modern Roman}
\usepackage{amsmath,amssymb,mathtools,unicode-math}
\setmathfont{Latin Modern Math}
\usepackage{enumitem,needspace,xurl,hyperref}
\setlength{\parindent}{0pt}
\setlength{\parskip}{5pt}
\setlength{\emergencystretch}{3em}
\newcommand{\N}{\mathbb{N}}
\newcommand{\Z}{\mathbb{Z}}
\newcommand{\Q}{\mathbb{Q}}
\newcommand{\R}{\mathbb{R}}
\newcommand{\C}{\mathbb{C}}
\newcommand{\F}{\mathbb{F}}
\newcommand{\A}{\mathbb{A}}
\newcommand{\PP}{\mathbb{P}}
\newcommand{\E}{\mathbb{E}}
\newcommand{\eps}{\varepsilon}
\newcommand{\dd}{\,\mathrm{d}}
\begin{document}
\section*{UnsolvedMath ID: 30003672}
\textbf{Source problem code:} catalogue entry
\subsection{Definitions and mathematical statement}
Let \(N_0\subset M_0\subset B(\mathcal H_0)\) be von Neumann algebras and let \(\Omega_0\) be a unit vector. Interpret a \emph{standard half-sided modular inclusion} as follows: \(\Omega_0\) is cyclic and separating for both \(N_0\) and \(M_0\), is cyclic for \(N_0'\cap M_0\), and the modular operator \(\Delta_{M_0,\Omega_0}\) satisfies
\[
\Delta_{M_0,\Omega_0}^{it}N_0\Delta_{M_0,\Omega_0}^{-it}\subseteq N_0
\qquad\text{whenever }\varepsilon t\geq 0,
\]
for a fixed sign \(\varepsilon\in\{+1,-1\}\). Here cyclicity means that the orbit of the vector under the algebra has dense linear span, and separating means that \(x\Omega_0=0\) implies \(x=0\) for elements of the algebra. The supplied record does not specify the half-sided sign convention; either convention must be used consistently throughout the construction.

Let \(K\) be a finite nonempty indexing set, and for each \(k\in K\) take a copy \((N_k\subset M_k,\Omega_k)\) of \((N_0\subset M_0,\Omega_0)\), with vector state \(\varphi_k(x)=\langle\Omega_k,x\Omega_k\rangle\). Form the reduced von Neumann algebra free product with respect to these states,
\[
(M,\varphi)=\mathop{*}_{k\in K}(M_k,\varphi_k),
\]
and let \(\Omega\) be its cyclic vector in the free-product representation. Write \(\lambda_k:M_k\to M\) for the canonical embeddings and set
\[
N=\left(\bigcup_{k\in K}\lambda_k(N_k)\right)''.
\]
The free-product state restricts to \(\varphi_k\) on each copy, and freeness means that
\[
\varphi\bigl(\lambda_{k_1}(x_1)\cdots\lambda_{k_r}(x_r)\bigr)=0
\]
whenever \(r\geq1\), \(x_j\in M_{k_j}\), \(\varphi_{k_j}(x_j)=0\), and \(k_j\ne k_{j+1}\) for adjacent indices.

\textbf{Problem.} For the supplied initial inclusion and finite family of copies, determine the relative commutant
\[
R=N'\cap M=\{x\in M: xn=nx\text{ for every }n\in N\}.
\]
In particular, determine whether \(R\ne\C1\). If it is nontrivial, identify the M\"obius-covariant net associated with the inclusion, insofar as the necessary reconstruction hypotheses hold. Such a net assigns von Neumann algebras to proper nonempty open intervals of the circle and has isotony, locality, and positive-energy M\"obius covariance with a vacuum vector.

\textbf{Unresolved qualifications.} The original wording calls \(K\) a finite group, whereas the reconstructed statement uses a finite set. The construction above uses only its underlying indexing set; no group action or use of multiplication in \(K\) is specified. Nonemptiness is imposed to describe a family of copies, since the source does not discuss an empty free product or give a lower bound on \(|K|\). Moreover, nontriviality of \(R\) alone does not state that \(\Omega\) is cyclic for \(R\), as required for standardness in the usual reconstruction from a standard half-sided modular inclusion. The record does not specify whether standardness of the resulting inclusion is part of the intended question, or whether reconstruction should instead occur on a suitable subspace. Thus the net-identification clause retains this unresolved qualification.
\subsection{Short English statement}
Take finitely many state-preserving copies of a standard half-sided modular inclusion and form their free-product inclusion. Determine which elements of the larger algebra commute with the smaller one. If this relative commutant is nontrivial, identify the associated M\"obius-covariant net, subject to clarifying the reconstruction hypotheses. The supplied record labels the problem open; this draft does not independently verify that status.
\subsection{Sources}
\begin{itemize}
\item Local UnsolvedMath registry, permanent ID \(30003672\), source problem code \texttt{catalogue entry}, in the supplied complete source record associated with \texttt{data/unsolvedmath/problems.json}; collection: Oberwolfach Reports Open Problems. No individual proposer is supplied; the recorded proposed year is 2017.
\item Supplied report citation: \emph{Reflection Positivity} (2018), Oberwolfach Reports, DOI \texttt{10.4171/owr/2017/55}, pp.\ 3263--3343. \url{https://doi.org/10.4171/owr/2017/55}. The supplied title year is 2018, while the report identifier and recorded proposed year refer to 2017; these details are retained as supplied.
\item Supplied literature source: \emph{Free products in AQFT}. \url{https://aif.centre-mersenne.org/item/10.5802/aif.3269.pdf}.
\item Supplied literature source: \emph{OWR Report 55/2017}. \url{https://ems.press/content/serial-article-files/46717}.
\item Supplied literature assessment, dated exactly \texttt{2026-08-21}: the relative commutant problem for finite free-product half-sided modular inclusions is described as still unknown, while nontrivial relative commutants are reported in related finite M\"obius-net constructions. This is supplied dated curation, not an independently verified literature assessment; related constructions are not asserted here to settle the stated problem.
\item The registry labels the statement \texttt{corrected\_verified} and reports that the official report was checked and that the reconstruction restored the free-product inclusion and finite indexing set. These are supplied provenance claims, not source verification performed for this draft.
\end{itemize}
\end{document}
